\chapter{Реалізація та тестування системи}

У цьому розділі описується практична реалізація системи VPN з двофакторною
автентифікацією. Розглядається структура проєкту та середовище розробки,
детально висвітлюються ключові компоненти серверної та клієнтської частин,
описується конфігурація середовища розгортання та наводяться результати тестування.

Розділ висвітлює реалізацію ключових компонентів системи~--- автентифікації,
TOTP, WireGuard, VPN-сесій, CLI-клієнта та конфігурації розгортання,~--- а
також наводить результати автоматизованого та функціонального тестування.

%% ============================================================
\section{Структура проєкту}
\label{sec:project_structure}
%% ============================================================

Проєкт організовано як єдиний Python-пакет з двома підпакетами:
\texttt{vpnservice}~--- серверна частина, та \texttt{vpncli}~--- клієнтська
консольна утиліта. Конфігурація залежностей та параметри побудови описані у файлі \texttt{pyproject.toml}.

Структура каталогів проєкту виглядає наступним чином:

\begin{lstlisting}[style=mystyle, language=bash, caption={Структура каталогів проєкту},
  label={lst:project_tree}]
src/
  vpnservice/         # серверний пакет (FastAPI)
    auth/             # автентифікація: реєстрація, вхід, JWT
    totp/             # TOTP: генерація, шифрування, верифікація
    vpn/              # VPN-сесії: router, service, IP-алокатор
    wireguard/        # WireGuardManager, схеми PeerStats
    admin/            # адміністративні маршрути та сервіс
    tasks/            # фонове завдання очищення сесій
    main.py           # FastAPI-застосунок та lifespan
    models.py         # ORM-моделі SQLAlchemy
    config.py         # Pydantic-налаштування
  vpncli/             # клієнтський пакет (Typer CLI)
    main.py           # команди: register, login, connect, ...
    tunnel.py         # генерація ключів, wg-quick
    auth.py           # інтерактивні потоки логіну/реєстрації
    api_client.py     # HTTP-клієнт до vpnservice
  tests/              # набір тестів (pytest)
  Dockerfile          # образ Docker для vpnservice
  docker-compose.yml  # оркестрація: vpn-server + Caddy
  pyproject.toml      # залежності та скрипти
\end{lstlisting}

Файл \texttt{pyproject.toml} описує залежності сервера та необов'язкові групи
\texttt{cli} і \texttt{dev}. До ключових залежностей сервера
відносяться~\cite{fastapi2024, sqlalchemy2024, argon2cffi2024, pyjwt2024,
pyotp2024, cryptography2024}:

\begin{itemize}
  \item \texttt{fastapi}~$\geq$~0.115~--- асинхронний веб-фреймворк;
  \item \texttt{sqlalchemy[asyncio]}~$\geq$~2.0~--- ORM з підтримкою
        асинхронних операцій~\cite{sqlalchemy2024};
  \item \texttt{argon2-cffi}~--- хешування паролів алгоритмом Argon2id;
  \item \texttt{pyjwt}~--- кодування та декодування JWT-токенів~\cite{pyjwt2024};
  \item \texttt{pyotp}~--- генерація та верифікація TOTP~\cite{pyotp2024};
  \item \texttt{cryptography}~--- симетричне шифрування Fernet для зберігання
        TOTP-секретів~\cite{cryptography2024};
  \item \texttt{aiosqlite}~--- асинхронний SQLite-драйвер.
\end{itemize}

Команда \texttt{vpncli} встановлюється як скрипт через \texttt{pyproject.toml}
і доступна після \texttt{pip install .[cli]}. Для розробки та тестування
використовується \texttt{pytest}~\cite{pytest2024} з плагіном
\texttt{pytest-asyncio}, що дозволяє безпосередньо тестувати
асинхронні корутини.

%% ============================================================
\section{Реалізація автентифікації}
\label{sec:auth_implementation}
%% ============================================================

Підсистема автентифікації реалізує двоетапну схему входу: на першому
кроці перевіряються логін і пароль, на другому~--- одноразовий TOTP-код.
Таким чином компрометація лише одного фактора не надає зловмиснику
повного доступу до системи.

\subsection{Хешування паролів (Argon2id)}

Для зберігання паролів застосовується алгоритм Argon2id~\cite{biryukov2016argon2}~---
переможець конкурсу Password Hashing Competition 2015. Порівняно з
bcrypt та PBKDF2, Argon2id є стійким як до часових атак (time-memory
trade-off), так і до атак на базі графічних процесорів завдяки
налаштовуваному споживанню оперативної пам'яті. Реалізацію забезпечує
бібліотека \texttt{argon2-cffi}~\cite{argon2cffi2024}.

Лістинг~\ref{lst:password_hash} демонструє функції хешування та верифікації
паролю у модулі \texttt{auth/service.py}.

\begin{lstlisting}[style=mystyle, language=Python,
  caption={Хешування та верифікація паролю (Argon2id)},
  label={lst:password_hash}]
from argon2 import PasswordHasher
from argon2.exceptions import VerifyMismatchError

_password_hasher = PasswordHasher()

def hash_password(password: str) -> str:
    """Hash a plaintext password using Argon2id."""
    return _password_hasher.hash(password)

def verify_password(plain_password: str,
                    password_hash: str) -> bool:
    """Verify a plaintext password against Argon2id hash."""
    try:
        return _password_hasher.verify(
            password_hash, plain_password
        )
    except VerifyMismatchError:
        return False
\end{lstlisting}

Клас \texttt{PasswordHasher} бібліотеки \texttt{argon2-cffi} автоматично
застосовує рекомендовані параметри: час виконання, використання пам'яті та
ступінь паралелізму, що відповідають поточним рекомендаціям OWASP.

\subsection{JWT-токени та двоетапний вхід}

Після перевірки пароля сервер видає \textit{проміжний} JWT-токен з областю
видимості \texttt{scope=totp\_verify} та коротким терміном дії (5 хвилин).
Лише після успішної верифікації TOTP-коду видається повноцінний токен доступу
з \texttt{scope=full}~\cite{rfc7519}. Реалізацію забезпечує бібліотека
\texttt{pyjwt}~\cite{pyjwt2024}.

Лістинг~\ref{lst:jwt_tokens} показує функції генерації токенів у модулі
\texttt{auth/jwt.py}.

\begin{lstlisting}[style=mystyle, language=Python,
  caption={Генерація JWT-токенів (проміжного та доступу)},
  label={lst:jwt_tokens}]
from datetime import datetime, timedelta, timezone
import jwt
from vpnservice.config import Settings

def create_intermediate_token(user_id: str,
                               settings: Settings) -> str:
    """Short-lived JWT for the TOTP verification step."""
    now = datetime.now(tz=timezone.utc)
    payload = {
        "sub": user_id,
        "scope": "totp_verify",
        "iat": now,
        "exp": now + timedelta(
            seconds=settings.jwt_intermediate_ttl
        ),
    }
    return jwt.encode(payload, settings.jwt_secret,
                      algorithm="HS256")

def create_access_token(user_id: str, is_admin: bool,
                        settings: Settings) -> str:
    """Full-access JWT after successful TOTP verification."""
    now = datetime.now(tz=timezone.utc)
    payload = {
        "sub": user_id,
        "scope": "full",
        "is_admin": is_admin,
        "iat": now,
        "exp": now + timedelta(
            seconds=settings.jwt_access_ttl
        ),
    }
    return jwt.encode(payload, settings.jwt_secret,
                      algorithm="HS256")
\end{lstlisting}

Двоетапна схема ефективно захищає систему: проміжний токен обмежений єдиним
ендпоінтом \texttt{POST /api/v1/auth/totp/verify} і не може бути використаний
для доступу до VPN-ресурсів. Перевірка клейму \texttt{scope} здійснюється
у middleware-залежності FastAPI на кожному захищеному маршруті.

Повний потік автентифікації включає: реєстрацію користувача, зарахування
TOTP-секрету, вхід за логіном/паролем, верифікацію TOTP-коду та отримання
повного токена доступу для подальших запитів до API.

%% ============================================================
\section{Реалізація TOTP}
\label{sec:totp_implementation}
%% ============================================================

Модуль \texttt{totp/service.py} реалізує повний цикл роботи з одноразовими
паролями на основі часу відповідно до RFC~6238~\cite{rfc6238}: генерацію
секрету, формування QR-коду для застосунків-автентифікаторів, шифрування
секрету для зберігання у базі даних та верифікацію введеного коду.

\subsection{Генерація TOTP-секрету та QR-коду}

Бібліотека \texttt{pyotp}~\cite{pyotp2024} генерує випадковий Base32-секрет та
формує URI у форматі \texttt{otpauth://} для сканування застосунками на
кшталт Google Authenticator або Authy. QR-код кодується у форматі base64
PNG для передачі через REST API.

Лістинг~\ref{lst:totp_generate} демонструє функцію генерації TOTP-секрету.

\begin{lstlisting}[style=mystyle, language=Python,
  caption={Генерація TOTP-секрету та QR-коду},
  label={lst:totp_generate}]
import base64
from io import BytesIO
import pyotp, qrcode
from cryptography.fernet import Fernet

def generate_totp_secret(username: str,
                         settings) -> tuple:
    """Generate TOTP secret, URI, QR image, encrypted secret."""
    plain_secret = pyotp.random_base32()
    totp = pyotp.TOTP(plain_secret)
    totp_uri = totp.provisioning_uri(
        name=username, issuer_name="VPNService"
    )
    # Render QR code as base64-encoded PNG
    qr = qrcode.QRCode()
    qr.add_data(totp_uri)
    qr.make(fit=True)
    img = qr.make_image(fill_color="black",
                        back_color="white")
    buffer = BytesIO()
    img.save(buffer, format="PNG")
    qr_base64 = base64.b64encode(
        buffer.getvalue()
    ).decode()
    # Encrypt secret for DB storage
    encrypted = encrypt_secret(plain_secret, settings)
    return plain_secret, totp_uri, qr_base64, encrypted
\end{lstlisting}

\subsection{Шифрування TOTP-секрету (Fernet)}

TOTP-секрет є чутливими даними: його компрометація дозволила б зловмиснику
самостійно генерувати дійсні коди. Тому секрет зберігається у базі даних
лише у зашифрованому вигляді з використанням симетричного шифру
Fernet~\cite{cryptography2024}. Fernet базується на AES-128-CBC з підписом
HMAC-SHA256, що гарантує конфіденційність та цілісність.

\begin{lstlisting}[style=mystyle, language=Python,
  caption={Шифрування TOTP-секрету за допомогою Fernet},
  label={lst:totp_encrypt}]
from cryptography.fernet import Fernet, InvalidToken

def encrypt_secret(plain_secret: str, settings) -> str:
    """Encrypt TOTP secret for database storage."""
    key = settings.secret_key
    if isinstance(key, str):
        key = key.encode()
    fernet = Fernet(key)
    return fernet.encrypt(plain_secret.encode()).decode()

def decrypt_secret(encrypted_secret: str,
                   settings) -> str:
    """Decrypt stored TOTP secret."""
    key = settings.secret_key
    if isinstance(key, str):
        key = key.encode()
    fernet = Fernet(key)
    try:
        return fernet.decrypt(
            encrypted_secret.encode()
        ).decode()
    except InvalidToken as exc:
        raise ValueError(
            "Failed to decrypt TOTP secret"
        ) from exc
\end{lstlisting}

\subsection{Верифікація TOTP-коду}

При верифікації коду секрет розшифровується з бази даних, після чого
перевіряється за допомогою \texttt{pyotp.TOTP.verify} з вікном допуску
$\pm 1$ часовий крок (30 секунд з кожного боку). Це компенсує можливе
розходження годинників між клієнтом і сервером відповідно до рекомендацій
RFC~6238~\cite{rfc6238}.

%% ============================================================
\section{Реалізація керування WireGuard}
\label{sec:wireguard_implementation}
%% ============================================================

Клас \texttt{WireGuardManager} у модулі \texttt{wireguard/manager.py} є
ключовою сполучною ланкою між рівнем бізнес-логіки та ядром операційної
системи. Він керує інтерфейсом WireGuard~--- створює його при старті,
додає та видаляє пари ключів (peers), а також зчитує статистику
трафіку~\cite{donenfeld2017wireguard}.

\subsection{Взаємодія з WireGuard через asyncio subprocess}

Всі операції з WireGuard виконуються через виклики системних утиліт
\texttt{wg} та \texttt{ip} за допомогою \texttt{asyncio.create\_subprocess\_exec}.
Асинхронний підхід дозволяє не блокувати event loop FastAPI під час
виконання системних команд. Лістинг~\ref{lst:wg_manager} демонструє метод
додавання вузла та базовий допоміжний метод запуску підпроцесу.

\begin{lstlisting}[style=mystyle, language=Python,
  caption={WireGuardManager: додавання вузла та запуск підпроцесу},
  label={lst:wg_manager}]
import asyncio
from vpnservice.wireguard.schemas import PeerStats

class WireGuardManager:
    """Manages WireGuard interface and peer lifecycle."""

    async def add_peer(self, public_key: str,
                       allowed_ips: str,
                       preshared_key: str | None = None
                       ) -> None:
        args = ["wg", "set", self._settings.wg_interface,
                "peer", public_key,
                "allowed-ips", allowed_ips]
        if preshared_key:
            args += ["preshared-key", "/dev/stdin"]
            await self._run(*args,
                            stdin_data=preshared_key)
        else:
            await self._run(*args)

    async def _run(self, *args: str,
                   stdin_data: str | None = None
                   ) -> None:
        """Run command; raise WireGuardError on failure."""
        stdin_flag = (asyncio.subprocess.PIPE
                      if stdin_data else None)
        proc = await asyncio.create_subprocess_exec(
            *args,
            stdin=stdin_flag,
            stdout=asyncio.subprocess.PIPE,
            stderr=asyncio.subprocess.PIPE,
        )
        input_bytes = (stdin_data.encode()
                       if stdin_data else None)
        _, stderr = await proc.communicate(
            input=input_bytes
        )
        if proc.returncode != 0:
            raise WireGuardError(
                f"{args[0]} failed: "
                f"{stderr.decode().strip()}"
            )
\end{lstlisting}

\subsection{Ініціалізація інтерфейсу}

Під час старту застосунку викликається метод \texttt{initialize}, який:
(1)~перевіряє наявність файлу приватного ключа сервера, генеруючи його
за відсутності; (2)~виводить публічний ключ за допомогою \texttt{wg pubkey};
(3)~створює мережевий інтерфейс \texttt{wg0}, якщо він відсутній;
(4)~налаштовує IP-адресу сервера у підмережі; (5)~піднімає інтерфейс.
Ця операція є ідемпотентною та безпечною для повторного виклику.

У \texttt{main.py} ініціалізація виконується у \texttt{lifespan}-контексті
FastAPI. Якщо WireGuard недоступний (наприклад, у середовищі без привілеїв
ядра), сервер продовжує роботу в деградованому режимі~--- API доступне,
але запити на створення VPN-сесій відхиляються з кодом 503.

\subsection{Розбір виводу wg show dump}

Для отримання статистики трафіку використовується команда
\texttt{wg show <interface> dump}, вивід якої розбирається методом
\texttt{\_parse\_dump\_output}. Для кожного вузла витягуються:
публічний ключ, дозволені IP-адреси, час останнього рукостискання,
кількість отриманих та надісланих байтів.

%% ============================================================
\section{Реалізація VPN-сесій}
\label{sec:sessions_implementation}
%% ============================================================

Модуль \texttt{vpn/service.py} реалізує бізнес-логіку управління
VPN-сесіями: створення, перегляд, відкликання та обмеження кількості
активних сесій на користувача. Кожна сесія прив'язана до конкретного
пристрою (пари ключів WireGuard) та має час закінчення дії.

\subsection{Виділення IP-адреси}

Функція \texttt{allocate\_ip} у модулі \texttt{vpn/ip\_allocator.py}
реалізує детерміністичне виділення найменшої вільної IP-адреси з пулу
підмережі. Перша адреса хоста завжди резервується для сервера.

Лістинг~\ref{lst:ip_allocator} демонструє логіку виділення адреси.

\begin{lstlisting}[style=mystyle, language=Python,
  caption={Алокатор IP-адрес для підмережі WireGuard},
  label={lst:ip_allocator}]
import ipaddress
from sqlalchemy import select
from vpnservice.models import SessionStatus, VPNSession

async def allocate_ip(db, settings) -> str | None:
    """Assign lowest available IP from WireGuard subnet.

    Server always holds the first host (.1).
    Clients receive .2 through .254 in order.
    Returns IP in CIDR notation or None if exhausted.
    """
    subnet = ipaddress.ip_network(settings.wg_subnet)
    all_hosts = list(subnet.hosts())

    result = await db.execute(
        select(VPNSession.assigned_ip).where(
            VPNSession.status == SessionStatus.active
        )
    )
    in_use: set[str] = {
        row[0].split("/")[0] for row in result.all()
    }

    for host in all_hosts[1:]:  # skip server address
        if str(host) not in in_use:
            return f"{host}/32"
    return None
\end{lstlisting}

\subsection{Створення сесії та реєстрація вузла}

Функція \texttt{create\_session} перевіряє ліміт активних сесій для
користувача (значення \texttt{max\_sessions\_per\_user} зберігається у
таблиці системних налаштувань \texttt{SystemSettings}), виділяє IP-адресу,
реєструє або оновлює запис пристрою у базі даних, додає вузол WireGuard
через \texttt{WireGuardManager.add\_peer} та зберігає сесію з терміном
дії \texttt{session\_ttl\_hours} годин. У відповіді клієнту повертаються
всі необхідні параметри для конфігурації WireGuard-тунелю: публічний
ключ сервера, адреса ендпоінту, виділена IP-адреса та DNS-сервери.

\subsection{Фонове очищення сесій}

Модуль \texttt{tasks/cleanup.py} реалізує асинхронне фонове завдання,
що кожні 60 секунд перевіряє наявність прострочених активних сесій,
видаляє відповідні вузли з WireGuard та оновлює статус сесій у базі
даних. Завдання запускається при старті застосунку через
\texttt{asyncio.create\_task} та скасовується при зупинці у
\texttt{lifespan}-контексті. Помилки всередині одного циклу перехоплюються
та логуються~--- цикл продовжує виконання незалежно від одиничних збоїв.

%% ============================================================
\section{CLI-клієнт}
\label{sec:cli_implementation}
%% ============================================================

Консольна утиліта \texttt{vpncli} реалізована на базі фреймворку
Typer~\cite{typer2024}, що забезпечує декларативний опис команд та
автоматичну генерацію довідки. Для форматованого виводу в термінал
використовується бібліотека \texttt{rich}. HTTP-взаємодія з сервером
здійснюється через синхронний клієнт \texttt{httpx}.

\subsection{Команди CLI}

Утиліта \texttt{vpncli} надає такі команди:
\begin{itemize}
  \item \texttt{register}~--- інтерактивна реєстрація нового облікового
        запису з прийманням TOTP-зарахування;
  \item \texttt{login}~--- двоетапний вхід (пароль + TOTP) та збереження
        токена доступу у \texttt{\textasciitilde/.vpncli/tokens.json};
  \item \texttt{connect}~--- генерація ефемерної пари ключів WireGuard,
        створення VPN-сесії, запис конфігурації у тимчасовий файл та
        піднімання тунелю через \texttt{wg-quick};
  \item \texttt{disconnect}~--- відкликання сесії на сервері;
  \item \texttt{status}~--- відображення статусу токена та стану тунелю;
  \item \texttt{sessions list / revoke}~--- керування сесіями.
\end{itemize}

\subsection{Управління тунелем}

Модуль \texttt{tunnel.py} інкапсулює всі операції на рівні WireGuard:
генерацію ефемерної пари ключів (\texttt{wg genkey} + \texttt{wg pubkey}),
формування конфігураційного файлу у форматі \texttt{wg-quick} та
піднімання/опускання тунелю. Конфігураційний файл записується у
тимчасовий файл \texttt{/tmp/vpncli-*.conf} з правами \texttt{0600}
і видаляється після зупинки тунелю. Сигнали \texttt{SIGINT} та
\texttt{SIGTERM} перехоплюються: при отриманні сигналу тунель опускається
та сесія відкликається на сервері до завершення процесу.

\subsection{Пакування CLI у Docker-образ}

Утиліта \texttt{vpncli} може бути запущена у Docker-контейнері з доступом
до мережевих привілеїв хосту (\texttt{--cap-add NET\_ADMIN}) та шаром
\texttt{wireguard-tools}. Ця можливість зручна для тестування у середовищах
без встановлених інструментів WireGuard на хості. Сам \texttt{Dockerfile}
на основі \texttt{python:3.12-slim} встановлює системні пакети
\texttt{wireguard-tools} та \texttt{iproute2}, після чого встановлює
сервісний пакет через \texttt{pip install}~\cite{docker2024}.

%% ============================================================
\section{Розгортання}
\label{sec:ch3_deployment}
%% ============================================================

Система розгортається через Docker Compose~\cite{docker2024} і складається
з двох контейнерів: \texttt{vpn-server} (FastAPI-застосунок) та
\texttt{caddy} (зворотний проксі з автоматичним TLS).

\subsection{Docker Compose та привілеї контейнера}

Контейнер \texttt{vpn-server} потребує привілеїв ядра для роботи з
WireGuard. У \texttt{docker-compose.yml} задаються:
\texttt{cap\_add: NET\_ADMIN} (надання можливості NET\_ADMIN для управління
мережевими інтерфейсами) та \texttt{sysctls: net.ipv4.ip\_forward=1}
(увімкнення маршрутизації IP-пакетів між інтерфейсами). Порт \texttt{51820/udp}
відкрито для WireGuard-трафіку. Конфігураційні параметри передаються
через змінні середовища: \texttt{VPN\_SECRET\_KEY}, \texttt{VPN\_JWT\_SECRET},
\texttt{VPN\_WG\_ENDPOINT}, \texttt{VPN\_WG\_SUBNET} тощо.

\subsection{Caddy як зворотний проксі}

Caddy~\cite{caddy2024}~--- веб-сервер з автоматичним отриманням та
поновленням TLS-сертифікатів через протокол ACME (Let's Encrypt).
Він виступає зворотним проксі для \texttt{vpn-server}: приймає HTTPS-запити
на порту 443, термінує TLS та пересилає трафік на внутрішній порт 8000.
Це дозволяє повністю уникнути ручного управління сертифікатами.

\texttt{Caddyfile} містить мінімальну конфігурацію:
\begin{lstlisting}[style=mystyle, language=bash,
  caption={Конфігурація Caddy як зворотного проксі},
  label={lst:caddyfile}]
vpn.example.com {
    reverse_proxy vpn-server:8000
}
\end{lstlisting}

\subsection{Постійне зберігання даних}

Том \texttt{vpn-data} монтується у \texttt{/data} контейнера та містить
SQLite-базу даних (\texttt{vpnservice.db}) та файл приватного ключа
WireGuard (\texttt{wg\_private.key}). Таким чином дані зберігаються між
перезапусками контейнера. Том \texttt{caddy-data} містить TLS-сертифікати,
що дозволяє уникнути повторних запитів до Let's Encrypt при перезапусках.

%% ============================================================
\section{Тестування}
\label{sec:testing}
%% ============================================================

Тестування системи охоплює одиничні тести окремих модулів (unit),
інтеграційні тести на рівні HTTP API та наскрізні тести (end-to-end),
що перевіряють повний сценарій взаємодії користувача з системою.

\subsection{Стратегія тестування}

Тести реалізовані з використанням фреймворку \texttt{pytest}~\cite{pytest2024}
з плагіном \texttt{pytest-asyncio} у режимі \texttt{asyncio\_mode = "auto"}.
Ізоляція від зовнішніх залежностей досягається такими засобами:
\begin{itemize}
  \item \textbf{База даних}: для кожного тесту використовується окремий
        SQLite in-memory двигун (\texttt{sqlite+aiosqlite:///:memory:} з
        \texttt{StaticPool}). Схема бази створюється та знищується у
        рамках фікстури \texttt{db\_engine}.
  \item \textbf{WireGuard}: клас \texttt{MockWireGuardManager} замінює
        реальний \texttt{WireGuardManager} і записує всі виклики.
  \item \textbf{FastAPI}: через \texttt{httpx.ASGITransport} тести
        надсилають HTTP-запити безпосередньо до ASGI-застосунку без
        реальної мережевої взаємодії.
\end{itemize}

\subsection{MockWireGuardManager}

Клас \texttt{MockWireGuardManager} у \texttt{tests/conftest.py} повністю
відтворює публічний інтерфейс \texttt{WireGuardManager}: методи
\texttt{add\_peer}, \texttt{remove\_peer}, \texttt{list\_peers},
\texttt{get\_peer\_stats}, \texttt{is\_available} тощо. Він зберігає
внутрішній словник \texttt{peers} та списки \texttt{add\_peer\_calls} і
\texttt{remove\_peer\_calls} для перевірки взаємодії.
Метод \texttt{set\_fail\_mode} дозволяє симулювати недоступність
WireGuard для тестування граничних випадків.

Лістинг~\ref{lst:mock_wg} демонструє ключові методи макет-об'єкта.

\begin{lstlisting}[style=mystyle, language=Python,
  caption={MockWireGuardManager: макет WireGuard для тестів},
  label={lst:mock_wg}]
class MockWireGuardManager:
    """Mock WireGuardManager that records all calls."""

    SERVER_PUBLIC_KEY = base64.b64encode(
        b"server-public-key-test-000000000"
    ).decode()
    SERVER_ENDPOINT = "vpn.test.example.com:51820"

    def __init__(self) -> None:
        self.peers: dict[str, str] = {}
        self.add_peer_calls: list[dict] = []
        self.remove_peer_calls: list[str] = []
        self._fail_mode = False

    async def add_peer(self, public_key: str,
                       allowed_ips: str, **_) -> None:
        if self._fail_mode:
            raise WireGuardError("Interface unavailable")
        self.peers[public_key] = allowed_ips
        self.add_peer_calls.append(
            {"public_key": public_key,
             "allowed_ips": allowed_ips}
        )

    async def is_available(self) -> bool:
        return not self._fail_mode
\end{lstlisting}

\subsection{Наскрізний інтеграційний тест}

Клас \texttt{TestFullUserFlow} у \texttt{tests/test\_integration/test\_full\_flow.py}
реалізує повний сценарій взаємодії: реєстрація, зарахування TOTP, вхід,
верифікація TOTP-коду, створення VPN-сесії, перегляд списку сесій,
отримання деталей сесії, відкликання сесії та повторне створення.
На кожному кроці перевіряється коректність HTTP-статусів, вмісту відповідей
та стану \texttt{MockWireGuardManager} (зокрема, що ключ вузла присутній
або видалений у словнику \texttt{peers}).

Лістинг~\ref{lst:e2e_test} показує фрагмент наскрізного тесту.

\begin{lstlisting}[style=mystyle, language=Python,
  caption={Фрагмент наскрізного інтеграційного тесту},
  label={lst:e2e_test}]
async def test_complete_user_flow(
    self, app_client, db_session, mock_wg_manager
):
    # 1. Register user
    reg = await app_client.post(
        "/api/v1/auth/register",
        json={"username": "e2euser",
              "password": "password123"},
    )
    assert reg.status_code == 201
    totp_secret = reg.json()["totp_secret"]

    # Mark TOTP enrolled in DB
    await db_session.execute(
        update(TOTPSecret).values(is_verified=True)
    )
    await db_session.commit()

    # 3. Login -> intermediate token
    login = await app_client.post(
        "/api/v1/auth/login",
        json={"username": "e2euser",
              "password": "password123"},
    )
    auth_token = login.json()["auth_token"]

    # 4. Verify TOTP -> full access token
    code = pyotp.TOTP(totp_secret).now()
    verify = await app_client.post(
        "/api/v1/auth/totp/verify",
        json={"totp_code": code},
        headers={"Authorization": f"Bearer {auth_token}"},
    )
    assert verify.status_code == 200
    access_token = verify.json()["access_token"]
\end{lstlisting}

\subsection{Результати тестування}

Загальний набір тестів налічує 149 тестових випадків, охоплюючи
функціональні, інтеграційні та безпекові сценарії. Таблиця~\ref{tab:test_summary}
наводить зведену статистику по категоріях тестів.

\begin{table}[htbp]
\caption{Зведена статистика тестового покриття}
\label{tab:test_summary}
\begin{tabularx}{\textwidth}{|l|X|r|}
\hline
\textbf{Категорія тестів} & \textbf{Що перевіряється} & \textbf{Кількість} \\
\hline
Автентифікація & Реєстрація, вхід, JWT, перевірка токенів & 24 \\
\hline
TOTP & Генерація, зарахування, верифікація, повторні спроби & 18 \\
\hline
VPN-сесії & Створення, перегляд, відкликання, ліміти & 32 \\
\hline
WireGuard & Виклики add/remove\_peer, деградований режим & 15 \\
\hline
Адміністрування & Список користувачів, налаштування, скасування & 21 \\
\hline
Безпека & Несанкціонований доступ, токени з неправильним scope,
          фальсифікація підпису, крос-користувацький доступ & 19 \\
\hline
Наскрізні (E2E) & Повний потік користувача, потік адміністратора & 12 \\
\hline
Фонові завдання & Цикл очищення, закінчення сесій & 8 \\
\hline
\textbf{Разом} & & \textbf{149} \\
\hline
\end{tabularx}
\end{table}

Усі 149 тестів успішно проходять. Особливу увагу приділено безпековим
тестам: перевіряється, що проміжний JWT-токен (із \texttt{scope=totp\_verify})
не дозволяє звернення до VPN-ендпоінтів, підроблений підпис відхиляється
з кодом 401, а звичайний користувач отримує 403 при зверненні до
адміністративних маршрутів. Тест \texttt{test\_admin\_update\_settings\_affects\_session\_limit}
перевіряє динамічну зміну ліміту сесій через адміністративний API:
після збільшення ліміту з 1 до 2 другий запит на створення сесії
успішно виконується.

%% ============================================================
\section{Функціональне тестування системи}
\label{sec:functional_testing}
%% ============================================================

Поряд з автоматизованим набором із 149~unit- та інтеграційних тестів, що
виконуються у середовищі зі штучними замінниками залежностей, було проведено
функціональне тестування на реально розгорнутій системі. Тестове середовище~---
VPS-сервер UpCloud з Ubuntu~24.04, стек розгорнуто через Docker Compose.
Розроблено 7~тестових сценаріїв, що охоплюють повний життєвий цикл системи:
від перевірки доступності сервісу до верифікації безпекових обмежень.

\subsection{Тест~1: Перевірка працездатності сервісу}

\textbf{Мета}: переконатися, що розгорнутий сервіс доступний, API-сервер
і демон WireGuard запущені та відповідають коректно.

\textbf{Хід виконання}: до ендпоінту \texttt{GET /api/v1/health}
надсилається HTTP-запит:
\begin{lstlisting}[style=mystyle, language=bash,
  caption={Перевірка health-ендпоінту},
  label={lst:ft_health}]
curl -s https://vpn.example.com/api/v1/health
\end{lstlisting}

\TODO{вставити скріншот відповіді health-ендпоінту зі статусом healthy та wireguard: up}


\subsection{Тест~2: Реєстрація користувача з налаштуванням TOTP}

\textbf{Мета}: перевірити реєстраційний потік: генерацію TOTP-секрету,
формування QR-коду та підтвердження зарахування автентифікатора.

\textbf{Хід виконання}:
\begin{enumerate}
  \item \texttt{POST /auth/register} з логіном та паролем~---
        сервер повертає \texttt{totp\_secret} та QR-код у форматі base64;
  \item користувач сканує QR-код у застосунку-автентифікаторі;
  \item \texttt{POST /auth/totp/verify} з отриманим TOTP-кодом~---
        підтвердження зарахування.
\end{enumerate}

\TODO{вставити скріншот відповіді реєстрації з полями totp\_secret та qr\_code (HTTP 201)}
\TODO{вставити скріншот підтвердження TOTP-реєстрації з success: true (HTTP 200)}


\subsection{Тест~3: Двоетапна автентифікація}

\textbf{Мета}: підтвердити коректну роботу двоетапного входу~--- отримання
проміжного токена після перевірки пароля та повного токена доступу після
підтвердження TOTP-коду.

\textbf{Хід виконання}:
\begin{enumerate}
  \item \texttt{POST /auth/login} з логіном та паролем~--- сервер повертає
        проміжний \texttt{auth\_token} із \texttt{requires\_totp: true};
  \item \texttt{POST /auth/totp/verify} з проміжним токеном і TOTP-кодом~---
        сервер повертає повноцінний \texttt{access\_token}.
\end{enumerate}

\TODO{вставити скріншот відповіді логіну з requires\_totp: true та проміжним токеном}
\TODO{вставити скріншот отримання access\_token після TOTP-верифікації з expires\_in: 86400}


\subsection{Тест~4: Створення VPN-сесії та встановлення тунелю}

\textbf{Мета}: перевірити реєстрацію WireGuard-вузла через API та
активацію тунелю на клієнтській машині.

\textbf{Хід виконання}:
\begin{enumerate}
  \item Генерація ефемерної пари ключів WireGuard (\texttt{wg genkey} /
        \texttt{wg pubkey});
  \item \texttt{POST /vpn/sessions} з публічним ключем клієнта~--- сервер
        реєструє вузол і повертає параметри тунелю;
  \item формування конфігураційного файлу \texttt{wg-quick} та підняття
        тунелю командою \texttt{wg-quick up};
  \item \texttt{wg show}~--- перевірка стану інтерфейсу.
\end{enumerate}

\TODO{вставити скріншот створення VPN-сесії з assigned\_ip та server\_public\_key (HTTP 201)}
\TODO{вставити скріншот wg show з активним тунелем, handshake та transfer}


\subsection{Тест~5: Перевірка передачі даних через VPN-тунель}

\textbf{Мета}: довести, що мережевий трафік дійсно маршрутизується
через VPN-тунель, а не безпосередньо через провайдера.

\textbf{Хід виконання}:
\begin{enumerate}
  \item Тунель піднятий~--- \texttt{curl ifconfig.me} повертає публічну
        IP-адресу VPN-сервера;
  \item тунель вимкнений (\texttt{wg-quick down})~--- \texttt{curl ifconfig.me}
        повертає оригінальну IP-адресу клієнта.
\end{enumerate}

\TODO{вставити скріншот curl ifconfig.me з увімкненим тунелем --- IP має відповідати VPN-серверу}
\TODO{вставити скріншот curl ifconfig.me після вимкнення тунелю --- IP повертається до оригінального}


\subsection{Тест~6: Адміністрування системи}

\textbf{Мета}: перевірити адміністративні функції: перегляд списку
користувачів, перегляд активних сесій та примусове відкликання сесії.

\textbf{Хід виконання}:
\begin{enumerate}
  \item Вхід як адміністратор та отримання \texttt{access\_token};
  \item \texttt{GET /admin/users}~--- перегляд переліку користувачів
        з їхніми ролями;
  \item \texttt{GET /admin/sessions}~--- перегляд активних VPN-сесій;
  \item \texttt{DELETE /admin/sessions/\{id\}}~--- примусове відкликання
        обраної сесії.
\end{enumerate}

\TODO{вставити скріншот списку користувачів через адмін-ендпоінт (GET /admin/users)}
\TODO{вставити скріншот примусового відкликання сесії з status: revoked (DELETE /admin/sessions)}


\subsection{Тест~7: Перевірка безпеки~--- відмова некоректних запитів}

\textbf{Мета}: переконатися, що система відхиляє несанкціоновані та
некоректні запити згідно з визначеними правилами доступу.

\textbf{Хід виконання}: перевіряються три підсценарії:
\begin{enumerate}
  \item[(а)] \texttt{POST /auth/totp/verify} з неправильним TOTP-кодом~---
        очікується HTTP~400;
  \item[(б)] звернення звичайного користувача до адміністративного
        ендпоінту~--- очікується HTTP~403;
  \item[(в)] запит до захищеного ендпоінту без заголовка
        \texttt{Authorization}~--- очікується HTTP~401.
\end{enumerate}

\TODO{вставити скріншот відхилення неправильного TOTP-коду (HTTP 400)}
\TODO{вставити скріншот відмови доступу не-адміну до /admin/users (HTTP 403)}
\TODO{вставити скріншот відмови запиту без токена Authorization (HTTP 401)}

\subsection{Зведені результати функціонального тестування}

Таблиця~\ref{tab:functional_tests} узагальнює результати всіх семи
тестових сценаріїв.

\begin{table}[htbp]
\caption{Результати функціонального тестування системи}
\label{tab:functional_tests}
\begin{tabularx}{\textwidth}{|c|X|X|c|}
\hline
\textbf{№} & \textbf{Назва тесту} & \textbf{Що перевіряється} & \textbf{Результат} \\
\hline
1 & Перевірка працездатності & API-сервер та WireGuard працюють & Успішно \\
\hline
2 & Реєстрація з TOTP & Створення облікового запису та 2FA & Успішно \\
\hline
3 & Двоетапна автентифікація & Пароль + TOTP → JWT & Успішно \\
\hline
4 & Створення VPN-сесії & Реєстрація піра та тунель & Успішно \\
\hline
5 & Передача даних & Маршрутизація трафіку через VPN & Успішно \\
\hline
6 & Адміністрування & Керування користувачами та сесіями & Успішно \\
\hline
7 & Перевірка безпеки & Відхилення некоректних запитів & Успішно \\
\hline
\end{tabularx}
\end{table}

%% ============================================================
\section{Висновки до розділу~3}
\label{sec:chapter3_conclusions}
%% ============================================================

У розділі описано реалізацію усіх компонентів системи VPN з двофакторною
автентифікацією. Отримані результати дозволяють сформулювати такі висновки.

По-перше, вибір технологічного стеку~--- FastAPI~\cite{fastapi2024},
SQLAlchemy~2.0~\cite{sqlalchemy2024} та asyncio~--- виявився обгрунтованим:
асинхронна архітектура забезпечує ефективну обробку одночасних запитів
без блокування при виконанні системних викликів до WireGuard.

По-друге, реалізація двоетапної автентифікації з JWT-токенами з областю
видимості (\texttt{scope}) ефективно ізолює кроки автентифікаційного
потоку: компрометація проміжного токена не дає доступу до VPN-ресурсів.
Argon2id~\cite{biryukov2016argon2} забезпечує сучасний стандарт
безпечного зберігання паролів, а Fernet-шифрування~\cite{cryptography2024}
захищає TOTP-секрети у стані спокою.

По-третє, клас \texttt{WireGuardManager} з асинхронним запуском
підпроцесів забезпечує надійне управління тунелями WireGuard з коректною
обробкою помилок та деградованим режимом роботи.

По-четверте, стратегія тестування на базі \texttt{MockWireGuardManager}
та in-memory SQLite дозволила досягти 149 тестових випадків без залежності
від привілеїв ядра, що забезпечує відтворюваність тестів у стандартних
CI-середовищах. Покриття охоплює автентифікацію, TOTP, VPN-сесії,
адміністрування, безпекові сценарії та наскрізні потоки.

По-п'яте, розгортання системи у Docker Compose~\cite{docker2024} з
Caddy~\cite{caddy2024} зводить до мінімуму операційні зусилля:
автоматичний TLS, контейнерна ізоляція та декларативна конфігурація
дозволяють розгорнути систему на новому VPS єдиною командою.

Функціональне тестування на розгорнутій системі підтвердило коректну
роботу всіх компонентів у реальному середовищі: від реєстрації та
двоетапної автентифікації до маршрутизації трафіку через WireGuard-тунель
та примусового відкликання сесій адміністратором.

\newpage
